\input{../style}

\title{{\Large Programmation sur Processeur Graphique -- GPGPU }\\
\vspace{-0.4em}
{\large TD 2 : allocation mémoire, kernel et erreurs}\\
{\small Centrale Nantes}\\
\small P.-E. Hladik, \small{\texttt{pierre-emmanuel.hladik@ec-nantes.fr}}\\
---\\
{\scriptsize  Version 1.0.0 (\today)}\\
}


\begin{document}

\maketitle

%\begin{table}[htdp]
%\caption{Suivi des modifications}
%\begin{center}
%\begin{tabular}{|c|c|c|p{8cm}|}
%\hline
%Date & Version & Auteur & Modifications\\
%\hline
%JJ/MM/12 & 0.3.X & P.-E. Hladik&  \\
%\hline
%\end{tabular}
%\end{center}
%\label{default}
%\end{table}%


%%%%%%%%%%%%%%%%%%%%%%%%%
  \section{Le b.a.-ba du CUDA}
%%%%%%%%%%%%%%%%%%%%%%%%%

Vous allez reprendre le code d'addition de vecteur pour vous exercer.

\begin{objectif}[]{}
 \begin{itemize}
	\item utiliser {\tt cudaMemcpy} et {\tt cudaMalloc}
	\item lancer un kernel
\end{itemize}
\end{objectif}

\begin{wtodo}[Addition vectorielle]{}
\begin{enumerate}
\item Récupérer le code source {\tt addvec.cu}.
\item Modifiez le pour qu'il réalise le travail demander, pour cela il ne faut modifier que les \texttt{/* FIXE ME */}.
\item Compiler, exécuter et observer.
\item Modifier le code (valeur de N ligne 3 dans le code initial) pour prendre en considération une taille quelconque des vecteurs (et non pas simplement un multiple de 512).
\end{enumerate}
\end{wtodo}


\begin{howto}[recopie de la mémoire]{}
La documentation de {\tt cudaMemcpy} est disponible sur

\url{https://docs.nvidia.com/cuda/cuda-runtime-api/group__CUDART__MEMORY.html}
\end{howto} 

\begin{howto}[Lancement d'un kernel]{}
Le lancement d'un kernel se fait comme un appel normal à une fonction en ajoutant des paramètres de configuration d'exécution entre {\tt {<}{<}{<}} et {\tt {>}{>}{>}} avant les arguments de la fonction. Le premier argument donne le nombre de blocs et le second le nombre de threads par bloc.
\end{howto} 

\newpage
%%%%%%%%%%%%%%%%%%%%%%%%%
  \section{Des bonnes pratiques : gestion des erreurs}
%%%%%%%%%%%%%%%%%%%%%%%%%

\begin{objectif}[]{}
 \begin{itemize}
	\item gestion des erreurs synchrones
\end{itemize}
\end{objectif}

\begin{wtodo}[Gestion des erreurs]{}

Toutes les fonctions de l'API CUDA renvoient un code d'erreur de type {\tt cudaError\_t}. Si l'appel ne retourne pas d'erreur cette valeur est égale à {\tt cudaSuccess}.

Le fichier {\tt error.h} fourni avec le sujet propose une macro (voir ci-dessous ce qu'est une macro) pour tester tous les appels aux fonctions CUDA. La macro permet de générer à la compilation le code pour ajouter un test sur la valeur de retour de la fonction (et donc sur le succès ou non de la fonction). Pour pouvoir utiliser la macro il suffit d'inclure {\tt error.h} dans votre fichier.

Utiliser cette macro pour tester le code d'erreur de retour pour chaque appel à l'API dans le code {\tt error.cu} (voir ci-dessous comment l'utiliser dans un code).
\end{wtodo}


\begin{howto}[C'est quoi une macro en C]{}
Une macro en C est une construction définie par le préprocesseur qui permet de remplacer du code avant la compilation. Les macros sont définies avec le mot-clé {\tt \#define} et servent généralement à simplifier le code, améliorer sa lisibilité ou définir des constantes. Elles ne sont pas des fonctions ou des variables, elles sont simplement remplacées directement dans le code source avant la compilation, par exemple, {\tt \#define PI 3.14159}. il est aussi possible de passer des arguments à une macro, par exemple {\tt \#define SQUARE(x) ((x) * (x))}.
\end{howto} 

\begin{howto}[Créer un affichage lisible]{}
La fonction {\tt cudaGetErrorString(cudaError\_t err)} convertit un code d'erreur en une chaîne de caractères, il peut être utilisé par exemple 
\begin{verbatim}
printf("%s in %s at line %d\n", cudaGetErrorString(err), __FILE__, __LINE__);
exit(EXIT_FAILURE);
\end{verbatim}
\end{howto} 

\begin{howto}[Passer un paramètre à une macro]{}
Une macro avec paramètres s'écrit {\tt \#define nom (<liste\_paramètres>) macro}, par exemple :
\begin{verbatim}
#define min(a, b)       ((a) < (b) ? (a) : (b))
\end{verbatim}
\end{howto} 

\begin{howto}[Ecriture d'une macro sur plusieurs lignes]{}
Il est possible d'écrire une macro sur plusieurs lignes en ajoutant {\tt \textbackslash} (backslash) à la fin de la ligne, par exemple : 
\begin{verbatim}
#define MACRO_EXPORT (target)                   \
  #ifdef _MSC_VER                               \
    #if defined TARGET_EXPORTS                  \
      #define TAGET_API __declspec( dllexport ) \
    #else                                       \
      #define TARGET_API __declspec( dllimport )\
  #endif // _LIB 
\end{verbatim}
\end{howto} 

\begin{howto}[Comment utiliser la macro {\tt CHECK\_ERROR} ]{}
La macro {\tt CHECK\_ERROR} prend en paramètre une fonction CUDA (et donc son retour d'erreur), par exemple pour vérifier le retour d'erreur de la fonction {\tt cudaFree(a);}, il faut écrire {\tt CHECK\_ERROR(cudaFree(a));}.
\end{howto} 


%%%%%%%%%%%%%%%%%%%%%%%%%
  \section{Des bonnes pratiques : gestion des erreurs (le retour)}
%%%%%%%%%%%%%%%%%%%%%%%%%

\begin{objectif}[]{}
 \begin{itemize}
	\item gestion des erreurs asynchrones
\end{itemize}
\end{objectif}

\begin{wtodo}[Gestion des erreurs asyncrhones]{}

L'utilisation de {\tt {<}{<}< ... {>}{>}>()} n'est pas un appel à l'API CUDA et ne renvoie pas de code d'erreur.
De plus il est asynchrone, ce qui signifie que le thread émet une requête pour exécuter le kernel sur le GPU, puis continue sans attendre que le kernel soit terminé.

Une erreur peut être capturée avec {\tt cudaGetLastError()}\footnote{\url{https://docs.nvidia.com/cuda/cuda-runtime-api/group\_\_CUDART\_\_ERROR.html\#group\_\_CUDART\_\_ERROR\_1g3529f94cb530a83a76613616782bd233}} après le lancement du noyau, mais il peut arriver qu'elle ne soit pas encore signalée au moment de l'appel (le kernel étant en cours d'exécution).

Il est possible de capturer l'erreur asynchrone en synchronisant explicitement l'exécution du noyau CUDA avec le thread, par exemple en utilisant {\tt cudaDeviceSynchronize}\footnote{\url{https://docs.nvidia.com/cuda/cuda-runtime-api/group\_\_CUDART\_\_DEVICE.html\#group\_\_CUDART\_\_DEVICE\_1g10e20b05a95f638a4071a655503df25d}}. Cependant, la synchronisation explicite affecte les performances et le comportement de l'exécution, il n'est donc pas recommandé de l'utiliser en production. Voir ci-dessous une utilisation possible de ces fonctions.

Ajouter les éléments nécessaire pour vérifier les erreurs de l'appel au kernel avec l'exemple fourni dans le fichier {\tt error2.cu}.

\end{wtodo}

\begin{howto}[Exemple pour capturer une erreur asynchrone]{}
Exemple d'utilisation de {\tt cudaGetLastError} avec {\tt cudaDeviceSynchronize} :
\begin{verbatim}
dev<<<...>>>(...);
ret = cudaGetLastError();
if (debug) ret = cudaDeviceSynchronize();
\end{verbatim}
Remarquez que {\tt cudaDeviceSynchronize} renvoie une erreur si l'une des tâches précédentes a échoué.
\end{howto} 


%%%%%%%%%%%%%%%%%%%%%%%%%
  \section{Bien gérer sa mémoire}
%%%%%%%%%%%%%%%%%%%%%%%%%

\begin{objectif}[]{}
 \begin{itemize}
	\item approfondissement sur l'allocation de la mémoire
\end{itemize}
\end{objectif}

\begin{wtodo}[Manipulation mémoire]{}

Le code source {\tt mesh.c} fournit des fonctions pour manipuler un objet 3D décrit au format OBJ. Une surface polygonale est décrite par un ensemble de sommets (vertices) et un ensemble de faces composés de ces sommets.

Un objet est décrit dans la structure {\tt object\_t} qui comprend quatre champs :
\begin{itemize}
	\item {\tt int nb\_vertices} : le nombre de sommets de l'objet
	\item {\tt vertex\_t *vertex\_t} : un pointeur vers un tableau de {\tt vertex\_t}
	\item {\tt int nb\_faces} : le nombre de faces
	\item {\tt face\_t * faces} : un pointeur vers un tableau de {\tt face\_t}
\end{itemize}
De même, {\tt vertex\_t} et {\tt face\_t} sont des structures contenant les informations de position des sommets et des faces.

Le code actuel permet d'importer un fichier OBJ et d'appliquer une transformation ({\tt flip\_and\_expand}) sur cet objet. La transformation inverse les faces et applique un facteur sur les coordonnées des sommets.
\medskip

Vous devez passer en CUDA ce code pour appeler la fonction {\tt flip\_and\_expand} sur le GPU. Ne cherchez pas à paralléliser cette fonction, mais simplement à l'appeler en la passant en {\tt \_{}\_global\_{}\_} dans un unique thread ({\tt flip\_and\_expand<{}<{}<{}1,1{}>{}>{}>}). {\bf La difficulté de cet exercice est l'utilisation de la mémoire.} N'oubliez pas de changer l'extension du fichier en {\tt .cu} pour pouvoir le compiler avec {\tt nvcc}.
\medskip

Si vous êtes en avance et vous vous ennuyez, vous pouvez commencer à paralléliser le traitement d'un mesh...
\end{wtodo}



%%%%%%%%%%%%%%%%%%%%%%%%%%
%  \section{Performance}
%%%%%%%%%%%%%%%%%%%%%%%%%%
%
%\begin{objectif}[]{}
% \begin{itemize}
% 	\item comparer les performances d'un code C et d'un code CUDA
%\end{itemize}
%\end{objectif}
%
%\begin{wtodo}[La première classe]{}
%\begin{enumerate}
%\item Récupérer le code {\tt addVecWithoutKernel.cu}, complier et exécuter.
%\item Ajouter dans le code de {\tt addVec.cu} des mesures de performance (temps de chargement de la mémoire et temps d'exécution) et afficher les. Ecrivez des macros  dans un fichier d'en-tête pour vous simplifiez l'écriture du code permettant de faire une mesure (par exemple avec une macro pour commencer la mesure du temps et une pour terminer).
%\item Compiler et exécuter les deux versions (C et CUDA) de l'addition vectorielle.
%\item Comparer les performances. A-t-on gagné quelque chose ?
%\end{enumerate}
%\end{wtodo}
%
%\begin{howto}[Création et horodatage d'un événemet CUDA]{}
%On pourrait utiliser les timers de l'OS sur le CPU, mais pendant que le noyau GPU fonctionne, il se peut que nous effectuions des calculs de manière asynchrone sur l'hôte. Pour résoudre cela on va utiliser des {\tt cudaEvent\_t} (événements CUDA).
%
%Un événement dans CUDA est un horodatage GPU qui est enregistré à un moment spécifié par l'utilisateur. L'API est relativement facile à utiliser et ne comporte que deux appels : la création d'un événement et son horodatage. Par exemple, au début d'une séquence de code, nous demandons au moteur d'exécution CUDA d'enregistrer l'heure actuelle. Nous le faisons en créant puis en enregistrant l'événement :
%\begin{verbatim}
%cudaEvent_t start;
%cudaEventCreate(&start);
%cudaEventRecord(start, 0);
%\end{verbatim}
%Le second argument de {\tt cudaEventRecord} restera un mystère, pour les curieux\footnote{\url{https://docs.nvidia.com/cuda/cuda-runtime-api/group__CUDART__EVENT.html}}.
%\end{howto} 
%
%\begin{howto}[Synchronisation]{}
%Malheureusement, il y a un problème de synchronisation. La partie la plus délicate de l'utilisation des événements est due au fait que certains appels dans CUDA C sont asynchrones. C'est excellent du point de vue des performances, car cela signifie que nous pouvons calculer quelque chose sur le GPU et le CPU en même temps, mais cela rend la mesure du temps délicate.
%
%Vous devez imaginer les appels à {\tt cudaEventRecord()} comme une instruction pour enregistrer l'heure actuelle placée dans la file d'attente du GPU. Par conséquent, notre événement ne sera pas réellement enregistré tant que le GPU n'aura pas tout terminé avant l'appel à cudaEventRecord(). C'est précisément ce que nous voulons pour que notre événement d'arrêt mesure l'heure correcte. Mais nous ne pouvons pas lire en toute sécurité la valeur de l'événement d'arrêt avant que le GPU n'ait terminé son travail précédent et enregistré l'événement d'arrêt. Heureusement, nous disposons d'un moyen de demander au CPU de se synchroniser sur un événement, la fonction API d'événement {\tt cudaEventSynchronize()} :
%\begin{verbatim}
%cudaEvent_t start, stop;
%cudaEventCreate(&start);
%cudaEventCreate(&stop);
%cudaEventRecord(start, 0);
%// do something
%cudaEventRecord(stop, 0);
%cudaEventSynchronize(stop);
%\end{verbatim}
%\end{howto} 
%
%\begin{howto}[Calcul du temps écoulé entre deux événements]{}
%Il est possible de connaître le temps en ms séparant deux évènement avec la méthodes {\tt cudaEventElapsedTime}. Par exemple :
%\begin{verbatim}
%float elapsedTime;
%cudaEventElapsedTime(&elapsedTime, start, stop);
%printf("Time %3.1f ms\n", elapsedTime);
%\end{verbatim}
%\end{howto} 

%%%%%%%%%%%%%%%%%%%%%%%%%%
%  \section{Connaître sa machine (épisode 2)}
%%%%%%%%%%%%%%%%%%%%%%%%%%
%
%\begin{ad}[]{}
% \begin{itemize}
%	Ceci est un exercice bonus si vous vous ennuyez.
%\end{itemize}
%\end{ad}
%
%\begin{objectif}[]{}
% \begin{itemize}
%	\item Afficher de nouvelles caractéristiques de sa machine
%\end{itemize}
%\end{objectif}
%
%\begin{wtodo}[Lecture des paramètres du GPU]{}
%\begin{enumerate}
%\item Compléter {\tt properties.cu} pour afficher le nombre maximum de threads par bloc, la taille maximum des dimensions d'un grid et d'un block.
%\end{enumerate}
%
%\end{wtodo}


\end{document}